\documentclass[11pt]{article}
\usepackage[a4paper,margin=28mm]{geometry}
\usepackage{amsmath,amssymb,amsthm,hyperref}
\newtheorem{proposition}{Proposition}
\title{Four equal dice with every sufficiently large admissible face count}
\author{Research certificate, 30 September 2026}
\date{}
\begin{document}
\maketitle

This note gives a small explicit certificate. Four18-sided fair dice
already appear in Ford, Grime, Harshbarger and Pollock,
\emph{Go First Dice for Five Players and Beyond} (2023), p.~87,
\href{https://doi.org/10.2478/rmm-2023-0004}{doi:10.2478/rmm-2023-0004}.
The word below was found independently; no novelty is claimed for the
existence of18-sided examples or the concatenation consequence. Write a die collection as a word, assigning increasing
integer labels to successive positions.

\begin{proposition}
Four permutation-fair equal-sided dice exist with $M$ faces each for every
multiple $M$ of $6$ with $M\ge12$.
\end{proposition}
\begin{proof}
Define three palindrome blocks
\[
 A=12033021,\qquad B=32011023,\qquad C=02133120.
\]
The words
\[
 W_{12}=ABCCBA,\qquad W_{18}=ABCCABBCA
\]
have respectively $12$ and $18$ occurrences of each of the four letters.
Their distinct-letter subsequence counts are as follows; each entry in
a column applies to every ordered pattern of that length.
\[
\begin{array}{c|rrrr}
 &1&2&3&4\\\hline
 W_{12}&12&72&288&864\\
 W_{18}&18&162&972&4374
\end{array}
\]
These integers can be checked directly by the following recurrence.
For each pattern $p_1\cdots p_k$, initialize $d_0=1$ and $d_j=0$ for
$1\le j\le k$. Read the word from left to right. At a letter $x$, update
$d_j\gets d_j+d_{j-1}$ whenever $p_j=x$, processing indices in decreasing
order. The final $d_k$ counts the desired subsequences. The companion
program and independent certificate check all $64$ nonempty patterns.
The displayed fourth-order counts are respectively $12^4/24$ and
$18^4/24$, proving fairness.

Concatenating equal-composition fair words preserves fairness: for a
fixed length-$k$ pattern, selecting its first $j$ letters from a block
with count $a$ and its remaining letters from one with count $b$ gives
$a^j b^{k-j}/(j!(k-j)!)$ occurrences. Summation over $j$ gives
$(a+b)^k/k!$.
Every integer $M/6\ge2$ is $2u+3v$ with nonnegative integers $u,v$.
Concatenating $u$ copies of $W_{12}$ and $v$ copies of $W_{18}$ proves
the assertion.
\end{proof}

In particular, the tempting stronger condition
$v_p(M)\ge\lfloor n/p\rfloor$ for equal-sided fair dice is false:
$n=4$, $M=18$, and $p=2$ give a counterexample. The example makes no
assertion about the exceptional count $M=6$.

\paragraph{Reproducibility.}
\path{search_palindrome_18.cpp} enumerates all triples of relabelled
palindrome blocks, retains the $192$ three-wise fair triples, and finds
three of their $24$ distinct fourth-order defect vectors summing to zero.
\path{four_dice_18_faces_root_check.json} records a separate Python
integer dynamic-programming verification of the resulting word.
\end{document}
